\begin{frame}{Optimal Substructure}
\begin{block}{최적화 문제에서의 정의}원문제의 최적해를 구성하는 관련 부분해도 그 subproblem의 최적해여야 한다.\end{block}
\begin{exampleblock}{Matrix Path}$(i,j)$까지의 최적 경로가 위에서 왔다면, 그 위 cell까지의 prefix도 최적이어야 한다.\end{exampleblock}
\begin{alertblock}{교체 논증}prefix가 최적이 아니라면 더 작은 prefix로 바꾸어 전체 경로를 개선할 수 있어 원래 경로의 최적성과 모순이다.\end{alertblock}
\end{frame}
\begin{frame}{Overlapping Subproblems}
\begin{block}{정의}재귀적으로 정의한 subproblem state가 여러 계산 경로에서 반복된다.\end{block}
\begin{columns}[T]\column{.48\textwidth}\begin{block}{Call tree}같은 state가 여러 node로 복제\end{block}
\column{.48\textwidth}\begin{block}{Dependency DAG}같은 state를 하나로 합쳐 저장·재사용\end{block}\end{columns}
\end{frame}
\begin{frame}{Quick Sort와 Matrix Path}
\small\centering\begin{tabular}{lll}\toprule
&Quick Sort&Matrix Path\\\midrule
subproblem&partition subarray&grid prefix\\
관계&주로 disjoint&ancestor state overlap\\
재사용&필요 없음&핵심\\
전형적 전략&divide-and-conquer&dynamic programming\\\bottomrule
\end{tabular}
\dptakeaway{차이는 재귀 문법이 아니라 subproblem graph의 구조다.}
\end{frame}
\begin{frame}{Optimal Substructure만으로 DP인가?}
\begin{itemize}
\item optimal substructure는 최적화 recurrence의 정당성을 설명하지만, 이것만으로 DP가 되지는 않는다.
\item DP의 계산 이득은 같은 state가 반복되는 overlap과 저장·재사용에서 나온다.
\item Fibonacci처럼 최적화가 아닌 DP에서는 optimal substructure 대신 정확한 state recurrence가 핵심이다.
\item Greedy도 optimal substructure를 가질 수 있지만 greedy-choice property가 추가로 필요하다.
\end{itemize}
\end{frame}
\begin{frame}{Greedy Counterexample: Local Cell만 보지 말자}
\begin{columns}[T]
\column{.48\textwidth}\centering
\textbf{local greedy}\par\smallskip
\begin{tikzpicture}[x=1cm,y=.82cm]
\foreach \r/\vals in {0/{1,2,100},1/{3,100,1},2/{1,1,1}}
  \foreach \v [count=\c from 0] in \vals {
    \pgfmathtruncatemacro{\chosen}{(\r==0)||(\c==2)}
    \ifnum\chosen=1\node[dp current]at(\c,-\r){\v};
    \else\node[dp cell]at(\c,-\r){\v};\fi}
\end{tikzpicture}
\[
1+2+100+1+1=105
\]
\column{.48\textwidth}\centering
\textbf{optimal path}\par\smallskip
\begin{tikzpicture}[x=1cm,y=.82cm]
\foreach \r/\vals in {0/{1,2,100},1/{3,100,1},2/{1,1,1}}
  \foreach \v [count=\c from 0] in \vals {
    \pgfmathtruncatemacro{\chosen}{(\c==0)||(\r==2)}
    \ifnum\chosen=1\node[dp answer]at(\c,-\r){\v};
    \else\node[dp cell]at(\c,-\r){\v};\fi}
\end{tikzpicture}
\[
1+3+1+1+1=7
\]
\end{columns}
\dptakeaway{가능한 predecessor의 최적값을 비교하는 DP와 현재 값 하나를 확정하는 greedy는 다르다.}
\vfill\muted{DP가 작동하는 이유를 정리했다. 이제 동일한 설계 template를 grid와 문자열 문제에 적용한다.}
\end{frame}
